\section{Introduction}
\label{sec:introduction}

Let \(\cP\) be the set of \(255\) points of
\(\operatorname{PG}(7,2)\).  We fix three block collections
\(\cB_1,\cB_2,\cB_3\), each consisting of \(66\,929\)
four-dimensional vector subspaces of \(\FF_2^8\).  They arise as three
constituents in the locked presentation studied in the companion
construction paper; their construction lineage is described in
\cite{KLW2018}.  The present paper starts from that finite triple and
studies a different layer: relations among its point-incidence columns.

Write
\[
 A_i:\ZZ[\cB_i]\longrightarrow \ZZ[\cP],
 \qquad
 B\longmapsto \sum_{x\in \operatorname{PG}(B)}x.
\]
Every \(A_i\) has rank \(255\).  We compare the rational relation module
\(K_i=\ker(A_i\otimes\QQ)\) with the submodule generated by the smallest
possible point trades.  Those trades have three blocks in each leg and
come from the two reguli in a quotient \(\operatorname{PG}(3,2)\).

The three configurations have the same parameters and a common symmetry
group
\[
 \GammaG\cong (C_{17}\rtimes C_4)\times C_3,
 \qquad |\GammaG|=204,
\]
but their relation modules are not isomorphic.  This is already visible
in every rational irreducible type.  Despite that separation, the same
local six-block geometry generates all rational relations in all three
cases.

\begin{theorem}[Main theorem]
\label{thm:main}
For the locked triple \((\cB_1,\cB_2,\cB_3)\), the following hold.
\begin{enumerate}[label=\textup{(\roman*)}]
  \item The rational \(\QQ\GammaG\)-modules \(K_1,K_2,K_3\) are pairwise
  nonisomorphic, and their multiplicities differ in each of the eight
  rational irreducible types.
  \item The \(\QQ\)-span of the minimum opposite-regulus trades is \(K_i\)
  for every \(i\).
  \item The associated incidence toric ideal has no nonzero element of
  degree below three.  Its primitive cubics are precisely the minimum
  opposite-regulus trades; every such cubic is indispensable.
  \item If \(L_i\) is the integral lattice generated by the minimum
  trades and
  \[
     \Theta_i=\ker_{\ZZ}(A_i)/L_i,
  \]
  then \(\Theta_i\) is finite and has zero \(2\)-, \(3\)-, and
  \(17\)-primary parts.
  \item The complete primary integral quotients vanish at all three selected
  good primes:
  \[
    \Theta_i[p^\infty]=0
    \qquad
    (i=1,2,3,\ p\in\{137,219\,097,3\,288\,036\,131\}).
  \]
  The apparent selected-row \(\rh/\cB_1\) defects at these primes
  disappear when the complete minimum-circuit equation maps are imposed.
  \item For the complete \(\rh/\cB_1\) minimum-circuit equation map over
  \(R=\ZZ[\zeta_{17}]\),
  \[
    \coker(F_{\rh,1}^{\mathrm{full}})
    \otimes_R R\left[\frac1{2\cdot3\cdot17}\right]=0.
  \]
  Hence this component has no prime-ideal support away from the rational
  primes \(2\), \(3\), and \(17\).
\end{enumerate}
\end{theorem}

The finite-prime assertions at \(137\), \(219\,097\), and
\(3\,288\,036\,131\) are complete for the frozen triple.  The last item
is stronger but narrower in a different direction: it is global away from
\(2\), \(3\), and \(17\), yet applies only to the complete
\(\rh/\cB_1\) component.  The corresponding global gates for the other
five large maps remain open, so no support theorem or global Smith
classification for the complete groups \(\Theta_i\) is asserted.
Likewise, our toric conclusion concerns the first nonzero graded layer; it
does not assert that all higher-degree moves are generated by cubics.

The proof has three interfaces.  Geometry identifies and counts the
minimum circuits.  Semisimple representation theory converts selected
circuits into exact component-rank problems and proves rational
generation.  Integral localization first proves saturation at the tested
primes and exposes why control-prime row selection cannot certify a
good-prime defect; an exact cyclotomic two-minor Fitting certificate then
closes the complete \(\rh/\cB_1\) component away from \(2\), \(3\), and
\(17\).  All numerical claims are accompanied by exact finite
certificates; the certificate contract is described in
\cref{app:certificates}.

\paragraph{Organization.}
\Cref{sec:locked} fixes notation and the rational characters.
\Cref{sec:modules} decomposes the relation modules.
\Cref{sec:circuits} proves the local circuit statement and gives the
exhaustive census.  \Cref{sec:rational-generation} proves rational
generation, and \cref{sec:toric} determines the cubic toric layer.
\Cref{sec:integral} treats the primes \(2,3,17\).
\Cref{sec:137} proves the complete primary classifications at
\(137\), \(219\,097\), and \(3\,288\,036\,131\).
\Cref{sec:global-rho} gives the global localized theorem for the single
complete \(\rh/\cB_1\) component.  Open integral and higher-degree
questions are separated in \cref{sec:discussion}.
